\documentclass[11pt]{article}
\usepackage[a4paper,margin=1in]{geometry}
\usepackage{amsmath,amssymb,bm}
\usepackage{graphicx}
\usepackage{float}
\usepackage{hyperref}
\usepackage{enumitem}
\graphicspath{{./figures/}}

\title{Shear wave propagagation physics and numerical implementations in the quasi-incompressible acousto-elastic soft tissue regime}
\author{Gianmarco Pinton}
\date{\today}

\begin{document}
\maketitle

\begin{abstract}

    Ultrasound-driven shear wave propagation in soft tissue underlie shear wave
    elastography (SWE) and emerging elastomechanical hypotheses of ultrasonic
    neuromodulation. Accurate modeling of radiation-force–induced tissue motion is
    therefore of central importance. However, general-purpose finite element
    elastodynamic solvers are often computationally expensive and unnecessarily
    broad for shear-dominant ultrasound applications. In soft biological tissue, the
    large separation between longitudinal and shear wave speeds ($\sim 10^3$ ratio)
    creates distinct temporal scales, enabling sequential acoustic-to-elastic
    modeling strategies. We derive a reduced viscoelastic elastodynamic formulation tailored to
    ultrasound-driven problems by systematically applying near-incompressibility,
    small-strain linearization, and solenoidal force projection to the full Navier
    equations. The resulting Kelvin--Voigt shear wave equation retains only the
    transverse dynamics relevant to radiation-force–induced motion while
    incorporating shear viscosity for attenuation modeling.     An explicit
    finite-difference time-domain (FDTD) implementation is developed, with
    second-order accuracy in space and time for the elastic component and stable
    treatment of viscous dissipation under classical CFL constraints.
     Validation studies demonstrate recovery of theoretical shear-wave speeds,
    convergence under grid refinement, and physically consistent attenuation
    behavior. The reduced framework provides a computationally efficient and
    physically transparent alternative to full multiphysics elasticity solvers,
    accommodates spatial heterogeneity in shear modulus, density, and viscosity,
    and facilitates integration with acoustic simulation tools and supporting both
    clinical SWE modeling and mechanistic investigation of ultrasound-induced
    neuromodulation. It's lightweight implementation and computational efficiency broaden applications to
     inverse problems, iterative methods, and machine learning data generation tasks.

    \end{abstract}

\section{Introduction}

Ultrasound-driven mechanical phenomena in soft tissue play a central role in
several biomedical applications, most prominently shear wave
elastography (SWE)~\cite{sarvazyan1998elasticity,barr2015elastography,gennisson2013ultrasound,sigrist2017ultrasound} and ultrasonic neuromodulation~\cite{blackmore2019ultrasound,jones2022neuromodulation,pinton2012braintherapy}. In SWE, acoustic radiation
force~\cite{nightingale2002radiation,palmeri2011radiation,nyborg1965acoustic,torr1984acoustic} generates transient elastic displacements whose propagation encodes the
underlying shear modulus of tissue. In ultrasonic neuromodulation, related
elastomechanical interactions—often framed through elastodynamic or
radiation-force–mediated hypotheses—have been proposed as potential mechanisms
linking acoustic exposure to neural activity~\cite{blackmore2019ultrasound,kubanek2018neuromodulation,yoo2022focused}. These application domains share a
common requirement: accurate and computationally efficient modeling of elastic
wave generation and propagation in soft solids under ultrasound excitation.

Extensive prior work has developed high-fidelity simulation frameworks for SWE
based on general-purpose finite element methods (FEM)~\cite{hughes2000finite,zienkiewicz2005finite}, including widely used
packages such as LS-DYNA and Abaqus~\cite{bercoff2004supersonic,capuzzi2018numerical}. These tools can represent complex
geometries, heterogeneous materials, nonlinear constitutive behavior, and
boundary interactions, making them highly versatile for broad classes of
elastodynamic problems. However, this generality comes at substantial
computational and algorithmic cost. For ultrasound-specific problems—where the
primary interest lies in small, radiation-force–induced shear displacements in
nearly incompressible soft tissue—the full generality of large-scale FEM is
often unnecessary and can obscure the essential physics governing wave
generation, propagation, and measurement.


A central challenge in ultrasound-driven elastodynamics is the strong
separation of physical scales between acoustic and shear phenomena. Acoustic
waves in soft tissue propagate with effective longitudinal speeds on the order
of $10^3$~m/s~\cite{szabo2014diagnostic,duck1990physical}, governed by the large bulk modulus, whereas shear waves travel
at speeds typically near $1$--$5$~m/s~\cite{sandrin2003transient,tanter2008quantitative}. This disparity produces temporal scales
that differ by approximately three orders of magnitude, despite spatial domains
that remain comparable. As a result, fully coupled multiphysics simulations that
resolve both acoustic and elastic dynamics simultaneously are computationally
inefficient and often unnecessary. Instead, practical modeling strategies
commonly treat the acoustic field and the resulting radiation-force–driven
elastic response on separate timescales, coupling them through intermediate
source terms: the acoustic pressure or intensity is computed first (e.g., with
fast nearfield and angular spectrum methods or full-wave codes), the radiation
force is derived from it, and that force is then supplied as a body force to a
shear wave solver~\cite{yang2018gpu}. Such acoustic–elastic coupling approaches, including those based
on established acoustic simulators (e.g., FIELD~II~\cite{jensen1996field}, Fullwave~\cite{pinton2020fullwave}), are
well documented in the SWE literature~\cite{bercoff2004supersonic,barr2015elastography}.
Full finite-element modeling of
ARF-induced shear wave propagation in tissue-mimicking media is covered by
established guidelines~\cite{palmeri2017fem}, which describe a sequential pipeline
(compute 3-D ARF push, apply to an FEM, extract motion) and recommend commercial
codes (e.g., LS-DYNA, Abaqus); that approach is general but computationally
heavy.

Numerical simulation is widely used to validate and improve shear wave
imaging techniques. Most simulations of shear waves generated by an applied
acoustic radiation force rely on the finite element (FEM) method. FEM models
have been developed for ARFI imaging~\cite{palmeri2005fem}, for the shear
wave response to radiation force in breast~\cite{nightingale2000fem} and
liver~\cite{palmeri2008liver}, for viscoelastic
propagation~\cite{lee2012viscoelastic}, for thermal effects in soft
tissue~\cite{palmeri2004thermal}, and for related applications such as lesion
characterization~\cite{fahey2005}, tracking and reflected-wave
effects~\cite{palmeri2006}, and transversely isotropic
media~\cite{rouze2013}. Alternatives to full-field FEM include the finite
difference method~\cite{orescanin2011fdtd,langdon2015}, the pseudospectral
method~\cite{qiang2015}, and the Green's function
approach~\cite{gennisson2003,bercoff2004supersonic,calle2005,chatelin2016}.
Green's function methods apply to both elastic~\cite{calle2005} and
viscoelastic media~\cite{bercoff2004supersonic} and are well suited when the
response is needed at only a few observation points; full-field methods (FEM,
FDTD, pseudospectral) instead compute the entire transient response over the
domain. The choice between these strategies therefore depends on whether one
needs the full field, on spatial heterogeneity, and on the type of source, i.e.
radiation force versus prescribed mechanical excitation.

Among finite-difference approaches, dedicated 3-D FDTD simulations of shear
waves for complex modulus imaging have been implemented~\cite{orescanin2011fdtd}
in the context of inverse problems. That work assumes incompressibility and
pure shear (normal stress components set to zero), yielding six first-order
velocity--stress equations in 3-D with Kelvin--Voigt viscoelasticity on a
staggered grid with PML boundaries; the source is a vibrating needle chosen
to minimize compressional energy. The present formulation, by contrast, uses a
displacement-based update with an explicit solenoidal projection of the body
force $\mathbf{f}_S$ via a scalar Poisson (CG) solve each step, so that
arbitrary (e.g., radiation-force) sources can be applied without prescribing a
needle geometry. Both are reduced shear-only forward solvers: the Orescanin
et al.\ approach targets algebraic Helmholtz inversion (AHI) for complex
modulus reconstruction from harmonic excitation, whereas the present solver
emphasizes integration with acoustic simulators and Kelvin--Voigt
attenuation in the shear equation.

Green's function methods offer a different trade-off. GPU-accelerated
implementations~\cite{yang2018gpu} provide fast evaluation of shear waves from
a prescribed radiation force in elastic and viscoelastic media within the same
sequential acousto-elastic workflow: the approach superposes contributions from
each source point and is well suited to homogeneous media and sparse
observation points. The present FDTD formulation instead advances the full
displacement field on a grid at each time step, so that the response is
available everywhere at once and the same framework can accommodate spatial
heterogeneity and more general boundary conditions. The solver
described here aligns with sequential, shear-focused strategies—minimal
reduced equation (solenoidal projection, Kelvin--Voigt shear only), explicit
FDTD, and matrix-free conjugate-gradient projection for the body force—yielding
a lightweight option for integration with acoustic simulators.

The objective of the present work is thus to derive and implement a reduced,
purpose-built formulation of shear elastodynamics tailored specifically to
ultrasound-driven problems in soft solids. Rather than pursuing full
general-purpose elasticity, we seek to isolate the minimal physical model
required to accurately reproduce radiation-force–induced shear displacements and
their propagation, while reducing computational complexity. The formulation
accommodates spatial heterogeneity in shear modulus, density, and viscosity. This
reduction enables tight integration with acoustic simulation tools and supports
efficient multiphysics workflows relevant to SWE, elastomechanical hypotheses
of ultrasonic neuromodulation, and iterative methods for inverse problems.

The remainder of this manuscript develops the governing physics, numerical
formulation, and validation of this reduced elastodynamic framework.

\section{Reduced Elastodynamic Formulation}
\label{sec:reduced_model}

This section derives the governing equations of motion for ultrasound-driven
shear dynamics in soft tissue, beginning from the full linear elastodynamic
continuum description and systematically applying the physical reductions
established below. The main assumptions are near-incompressibility of soft
tissue ($K \gg \mu$), small-strain linearization, and timescale separation
between acoustic and shear dynamics; they are invoked explicitly in
Sec.~\ref{sec:near_incompressible} and the incompressible limit
(Sec.~\ref{sec:incompressible_limit}).
\label{sec:assumptions}

\subsection{Near-Incompressibility in Soft Tissue and Modulus Hierarchy}
\label{sec:near_incompressible}

A key simplification exploited in ultrasound-driven elastodynamics is the extreme
separation between volumetric (compressional) stiffness and shear stiffness in
soft tissue~\cite{fung1993biomechanics,duck1990physical}. In isotropic linear elasticity~\cite{auld1973acoustic,graff1991wave,achenbach1973wave}, the longitudinal and shear wave
speeds satisfy
\begin{align}
c_L^2 &= \frac{K+\frac{4}{3}\mu}{\rho},\\
c_S^2 &= \frac{\mu}{\rho},
\end{align}
where $\rho$ is density, $K$ is bulk modulus, and $\mu$ is shear modulus. Using
typical ultrasound values in soft tissue~\cite{duck1990physical} ($c_L \approx 1540$~m/s and
$\rho \approx 1000$~kg/m$^3$) gives
\begin{equation}
K \approx \rho\,c_L^2 \approx 2.4\times 10^{9}\ \text{Pa},
\end{equation}
where the correction $\frac{4}{3}\mu$ is negligible because $\mu \ll K$.
Meanwhile, typical shear-wave speeds in elastography~\cite{sandrin2003transient,barr2015elastography} ($c_S \approx 1$--$5$~m/s)
imply
\begin{equation}
\mu = \rho\,c_S^2 \approx 1\times 10^{3}\text{--}2.5\times 10^{4}\ \text{Pa}.
\end{equation}
Thus, in ultrasound-relevant soft tissue,
\begin{equation}
\frac{K}{\mu} \sim 10^{5}\text{--}10^{6},
\end{equation}
which is a far stronger hierarchy than one would infer from moderate values of
Poisson's ratio alone.

For completeness, the elastic constants satisfy
\begin{equation}
E = 2\mu(1+\nu), \qquad
\nu = \frac{3K-2\mu}{2(3K+\mu)}, \qquad
K = \frac{E}{3(1-2\nu)}.
\end{equation}
Under the hierarchy $K \gg \mu$, these reduce to the familiar incompressible
limits
\begin{equation}
\nu \approx \frac{1}{2}-\frac{\mu}{6K}, \qquad
E \approx 3\mu,
\end{equation}
so that the Young's modulus remains in the kPa range even though the bulk
modulus is in the GPa range. Numerically, combining the representative values
above yields $\nu = 0.499999\ \text{(approximately)}$, emphasizing that
``near-incompressible'' in the ultrasound sense corresponds to $\nu$ extremely
close to $0.5$ when $K$ is inferred from $c_L$.

\paragraph{Interpretation for ultrasound-driven shear wave propagation.}
The hierarchy $K \gg \mu$ implies that compressional disturbances propagate
rapidly and are energetically expensive in volumetric strain, whereas the
radiation-force–induced response of interest in elastography~\cite{sandrin2003transient,tanter2008quantitative} (and potentially in
elastomechanical hypotheses of neuromodulation~\cite{blackmore2019ultrasound,pinton2012braintherapy}) is dominated by slower shear
motion governed primarily by $\mu$. This motivates (i) shear-focused elastic
models and (ii) sequential acousto--elastic workflows where the acoustic field
is used to compute forcing, and the subsequent shear dynamics are simulated on a
much slower time axis.

In ultrasound-driven problems, the external body force arises from acoustic
radiation force. Because the acoustic field evolves on timescales three orders
of magnitude faster than shear propagation (Sec.~\ref{sec:assumptions}), the
force distribution may be computed independently using an acoustic simulation
and treated as a prescribed source term in Eq.~\eqref{eq:shear_reduced_kv}. The
resulting formulation constitutes a sequential acoustic-to-elastic coupling
appropriate for shear wave elastography and related elastomechanical
applications in ultrasonic neuromodulation.


\subsection{Full linear viscoelastodynamics}

Let $\mathbf{u}(\mathbf{x},t)$ denote the displacement field in an isotropic,
linearly viscoelastic continuum of density $\rho$. We adopt a Kelvin--Voigt
constitutive law~\cite{catheline2004viscoelastic,ferry1980viscoelasticity}, in which the Cauchy stress is decomposed into an elastic and
a viscous contribution,
\begin{equation}
\boldsymbol{\sigma}
=
\lambda\,\mathrm{tr}(\boldsymbol{\varepsilon})\mathbf{I}
+ 2\mu\,\boldsymbol{\varepsilon}
+ 2\eta\,\dot{\boldsymbol{\varepsilon}},
\label{eq:kv_stress}
\end{equation}
where $\boldsymbol{\varepsilon} = \tfrac{1}{2}\left(\nabla \mathbf{u} + (\nabla \mathbf{u})^T\right)$
is the small-strain tensor, $(\cdot)\dot{}$ denotes a time derivative, $\lambda$ and
$\mu$ are the Lam\'e parameters, and $\eta$ is the shear viscosity. Conservation of
linear momentum with body force density $\mathbf{f}$ yields
\begin{equation}
\rho \frac{\partial^2 \mathbf{u}}{\partial t^2}
=
\nabla\cdot \boldsymbol{\sigma} + \mathbf{f}.
\label{eq:momentum_balance}
\end{equation}
Substituting Eq.~\eqref{eq:kv_stress} into Eq.~\eqref{eq:momentum_balance} gives
\begin{equation}
\rho \frac{\partial^2 \mathbf{u}}{\partial t^2}
=
(\lambda+\mu)\nabla(\nabla\cdot\mathbf{u})
+ \mu \nabla^2\mathbf{u}
+ \eta \nabla^2 \frac{\partial \mathbf{u}}{\partial t}
+ \mathbf{f},
\label{eq:navier_kv}
\end{equation}
where the viscous term contributes a Laplacian of particle velocity.


\subsection{Helmholtz decomposition of displacement}

The displacement field may be decomposed into irrotational (longitudinal) and
solenoidal (shear) components~\cite{auld1973acoustic,graff1991wave},

\begin{equation}
\mathbf{u} = \nabla \phi + \nabla \times \mathbf{\Psi},
\end{equation}

where $\phi$ is a scalar potential and $\mathbf{\Psi}$ is a vector potential
satisfying $\nabla \cdot \mathbf{\Psi} = 0$.

Substituting into Eq.~\eqref{eq:navier_kv} yields decoupled wave equations for the
longitudinal and shear components:

\begin{align}
\frac{\partial^2 \phi}{\partial t^2} &= c_L^2 \nabla^2 \phi + \frac{1}{\rho} f_L,
\\
\frac{\partial^2 \mathbf{\Psi}}{\partial t^2} &= c_S^2 \nabla^2 \mathbf{\Psi} + \frac{1}{\rho} \mathbf{f}_S,
\end{align}

where $f_L$ and $\mathbf{f}_S$ denote the longitudinal and solenoidal
projections of the body force.

\subsection{Incompressible limit and reduced shear equation with viscosity}
\label{sec:incompressible_limit}

Given the extreme modulus hierarchy $K \gg \mu$ (Sec.~\ref{sec:near_incompressible}),
the incompressible limit is taken:
\begin{equation}
\nabla \cdot \mathbf{u} = 0,
\label{eq:incompressible}
\end{equation}
which eliminates longitudinal dynamics and introduces a pressure field $p$
as a Lagrange multiplier that enforces the constraint. In this limit,
Eq.~\eqref{eq:navier_kv} becomes
\begin{equation}
\rho \frac{\partial^2 \mathbf{u}}{\partial t^2}
=
-\nabla p + \mu \nabla^2 \mathbf{u}
+ \eta \nabla^2 \frac{\partial \mathbf{u}}{\partial t}
+ \mathbf{f}.
\label{eq:incompressible_kv}
\end{equation}
Taking the divergence of this momentum equation and enforcing
Eq.~\eqref{eq:incompressible} yields a Poisson equation for $p$ that
keeps the displacement divergence-free at each time. Projecting
Eq.~\eqref{eq:incompressible_kv} onto the divergence-free subspace removes
the pressure gradient and gives the reduced forced shear wave equation
\begin{equation}
\rho \frac{\partial^2 \mathbf{u}}{\partial t^2}
=
\mu \nabla^2 \mathbf{u}
+ \eta \nabla^2 \frac{\partial \mathbf{u}}{\partial t}
+ \mathbf{f}_S,
\label{eq:shear_reduced_kv}
\end{equation}
where $\mathbf{f}_S$ is the solenoidal projection of the applied body force.
Equation~\eqref{eq:shear_reduced_kv} governs shear wave propagation with
wavespeed $c_S = \sqrt{\mu/\rho}$ and frequency-dependent attenuation set by $\eta$.
Under the assumptions of linear elasticity, near incompressibility, small
displacement, and timescale separation between acoustic and shear dynamics,
this equation forms the physical basis for the numerical formulation of
ultrasound-driven tissue motion.




\section{Numerical Methods}
\subsection{FDTD Time-Domain Update (Explicit, Second Order)}
Let $\bm{u}^n$ denote displacement at time $t_n=n\Delta t$. The implemented update~\cite{taflove2005computational} is
\begin{equation}
\bm{u}^{n+1}
=
2\bm{u}^{n} - \bm{u}^{n-1}
+
\Delta t^2 \, \bm{a}^{n},
\end{equation}
with acceleration
\begin{equation}
\bm{a}^{n}
=
\frac{\mu}{\rho}\nabla_h^2 \bm{u}^{n}
+
\frac{\eta}{\rho}\nabla_h^2 \dot{\bm{u}}^{n}
+
\frac{1}{\rho}\bm{f}_S^{n},
\quad
\dot{\bm{u}}^{n}\approx\frac{\bm{u}^{n}-\bm{u}^{n-1}}{\Delta t}.
\end{equation}
$\nabla_h^2$ is a second-order finite-difference 3D Laplacian on cell centers.

For homogeneous media without viscosity, a standard stability bound is approximately
\begin{equation}
\Delta t \lesssim \frac{1}{c_s}
\left(
\frac{1}{\Delta x^2}+\frac{1}{\Delta y^2}+\frac{1}{\Delta z^2}
\right)^{-1/2},
\end{equation}
which motivates the CFL-based~\cite{courant1928partial} step choice in validation scripts. The overall scheme is second-order accurate in time for the elastic component
and first-order accurate in time for viscous dissipation, consistent with
Kelvin--Voigt damping~\cite{catheline2004viscoelastic}. Numerical convergence and stability are verified empirically in
Sec.~\ref{sec:results} through grid-refinement and timestep studies.

\subsection{CG Projection for the Scalar Poisson Problem}
\subsubsection{Helmholtz Projection of Body Force}
At each step, the code projects $\bm{f}$ into solenoidal and irrotational parts:
\begin{equation}
\bm{f} = \bm{f}_S + \nabla \phi,
\quad
\nabla\cdot\bm{f}_S=0.
\end{equation}
Taking divergence gives Poisson equation:
\begin{equation}
\nabla^2 \phi = \nabla\cdot\bm{f},
\end{equation}
then
\begin{equation}
\bm{f}_S = \bm{f} - \nabla\phi.
\end{equation}
Numerically, this is solved by conjugate-gradient~\cite{golub2013matrix,shewchuk1994conjugate} on the scalar Poisson problem.

\paragraph{Physical consequence.}
The solver removes forcing components that would excite compressional motion. It is a shear-mode propagator, not a full elastic-wave solver.

\subsubsection{Why Conjugate Gradient on a Scalar Poisson Problem?}
The projection step enforces $\nabla\cdot\bm{f}_S=0$ by writing
\begin{equation}
\bm{f}_S = \bm{f} - \nabla\phi,
\end{equation}
which leads to the scalar Poisson equation
\begin{equation}
\nabla^2\phi = \nabla\cdot\bm{f}.
\end{equation}
Thus, only one scalar field $\phi$ must be solved each step (rather than a full vector system). After finite-difference discretization, this becomes a large sparse linear system
\begin{equation}
A\phi=b,
\end{equation}
where $A$ is the discrete negative Laplacian on interior nodes (with masking/boundary treatment).

Conjugate gradient (CG)~\cite{shewchuk1994conjugate,golub2013matrix} is a natural choice because:
\begin{itemize}[leftmargin=1.5em]
\item $A$ is symmetric positive-definite on the interior unknowns, matching CG requirements;
\item the 3D problem is large and sparse, so direct factorization is often memory-expensive;
\item CG only needs repeated stencil-like matrix-vector operations, which scale well in 3D.
\end{itemize}

From a physics perspective, this solve removes the irrotational (gradient) force component that would drive compressional response in full elastodynamics. The remaining divergence-free force preferentially excites transverse shear dynamics, consistent with the shear-focused SWE objective.

\subsubsection{Numerical Method Details for the Scalar Poisson CG Solve}
In implementation, the projection solve is performed in matrix-free form as
\begin{equation}
A\phi=b, \qquad A=-\nabla_h^2, \qquad b=\nabla_h\cdot\bm{f},
\end{equation}
where $\nabla_h^2$ is a second-order central-difference 3D Laplacian (7-point stencil on interior nodes).

\paragraph{Masked interior update.}
The code uses an interior mask $m$ (1 in interior, 0 on outer layer) and applies it after each update:
\begin{equation}
\phi \leftarrow \phi \odot m,\quad r \leftarrow r \odot m,\quad p \leftarrow p \odot m.
\end{equation}
This enforces the chosen boundary treatment and avoids boundary unknown drift.

\paragraph{CG iteration (matrix-free).}
Starting from $\phi_0=0$:
\begin{align}
r_0 &= (b-A\phi_0)\odot m, \qquad p_0=r_0,\\
\alpha_k &= \frac{r_k^T r_k}{p_k^T A p_k},\\
\phi_{k+1} &= (\phi_k + \alpha_k p_k)\odot m,\\
r_{k+1} &= (r_k - \alpha_k A p_k)\odot m,\\
\beta_k &= \frac{r_{k+1}^T r_{k+1}}{r_k^T r_k},\\
p_{k+1} &= (r_{k+1} + \beta_k p_k)\odot m.
\end{align}
No sparse matrix is assembled explicitly; $Aq$ is computed by applying the Laplacian stencil to $q$.

\paragraph{Stopping and numerical safeguards.}
The solver terminates when relative residual drops below tolerance,
\begin{equation}
\sqrt{\frac{\|r_k\|_2^2}{\|r_0\|_2^2}} < \text{tol},
\end{equation}
or when maximum iterations are reached. Additional guards stop iterations if denominator or residual becomes non-finite or non-positive in a way that indicates numerical breakdown.

\paragraph{Precision and cost.}
The Poisson CG routine is evaluated in float64 for improved robustness, then cast back to the field dtype. Computational cost is approximately $O(N\,n_{\mathrm{iter}})$ with $N$ grid points, which is efficient for repeated projection steps in 3D without the memory burden of direct factorization.


\subsection{Pre-computed Helmholtz Projection for Separable Body Forces}
\label{sec:precomputed_projection}

The CG projection described in Sec.~3.2 is applied at every time step in the
general solver, which is correct for arbitrary time-varying body forces
$\bm{f}(\bm{x},t)$. However, in the practically important case of acoustic
radiation force (ARF) driven shear wave excitation, the spatial and temporal
dependences of the body force are separable. The ARF spatial distribution
$\bm{b}_0(\bm{x})$ is determined entirely by the acoustic intensity field and
tissue absorption, while the temporal dependence is controlled by the push
envelope (e.g., a gated window during which the acoustic source is active). The
body force then takes the form
\begin{equation}
\bm{f}(\bm{x},t) = \bm{b}_0(\bm{x})\,g(t),
\label{eq:separable_force}
\end{equation}
where $g(t)$ is a scalar temporal envelope (for example, a step function that
equals one during the push and zero afterwards).

\subsubsection{Linearity of the Helmholtz Projection}

The Helmholtz--Hodge decomposition is a linear operation: for any vector field
$\bm{f}$ and scalar $\alpha$, the solenoidal projection $\mathcal{P}_S$
satisfies
\begin{equation}
\mathcal{P}_S[\alpha\,\bm{f}] = \alpha\,\mathcal{P}_S[\bm{f}].
\end{equation}
This follows directly from the linearity of the Poisson equation that defines
the projection. Given the decomposition $\bm{f} = \bm{f}_S + \nabla\phi$ with
$\nabla^2\phi = \nabla\cdot\bm{f}$, scaling the input by $\alpha$ scales both
$\phi$ and $\bm{f}_S$ by the same factor.

Applying this to the separable force in Eq.~\eqref{eq:separable_force}, the
solenoidal projection at any time $t_n$ is
\begin{equation}
\bm{f}_S^{\,n}
= \mathcal{P}_S\!\bigl[\bm{b}_0\,g(t_n)\bigr]
= g(t_n)\,\mathcal{P}_S[\bm{b}_0]
= g(t_n)\,\bm{f}_{S,0},
\label{eq:projected_separable}
\end{equation}
where $\bm{f}_{S,0} \equiv \mathcal{P}_S[\bm{b}_0]$ is the divergence-free
projection of the static spatial pattern, computed once.

\subsubsection{Algorithm}

The pre-computed projection replaces the per-step CG solve with a single
projection followed by scalar multiplication at each time step:

\begin{enumerate}[leftmargin=1.5em]
\item \textbf{Pre-computation (once).} Solve the scalar Poisson problem
\begin{equation}
\nabla_h^2 \phi_0 = \nabla_h \cdot \bm{b}_0,
\end{equation}
via conjugate gradient (as in Sec.~3.2), then form
\begin{equation}
\bm{f}_{S,0} = \bm{b}_0 - \nabla_h \phi_0.
\end{equation}

\item \textbf{Time stepping.} At each step $n$, the solenoidal force is
obtained by scaling:
\begin{equation}
\bm{f}_S^{\,n} = g(t_n)\,\bm{f}_{S,0}.
\end{equation}
The FDTD update (Sec.~3.1) then proceeds with the acceleration
\begin{equation}
\bm{a}^{n}
=
\frac{\mu}{\rho}\nabla_h^2 \bm{u}^{n}
+
\frac{\eta}{\rho}\nabla_h^2 \dot{\bm{u}}^{n}
+
\frac{g(t_n)}{\rho}\,\bm{f}_{S,0}.
\end{equation}
\end{enumerate}

\subsubsection{Computational Cost}

For a grid with $N$ points and $N_t$ time steps, the per-step projection
requires $N_t$ CG solves at cost $O(N \cdot n_{\mathrm{iter}})$ each, giving a
total projection cost of $O(N \cdot n_{\mathrm{iter}} \cdot N_t)$. The
pre-computed approach requires a single CG solve at cost
$O(N \cdot n_{\mathrm{iter}})$, plus $O(N)$ scalar--vector multiplication per
step, reducing the total projection cost to
$O(N \cdot n_{\mathrm{iter}} + N \cdot N_t)$.

In the ARF-driven regime the CG projection typically dominates wall-clock time
(the Laplacian stencil evaluations in the FDTD update are comparatively cheap).
For practical problem sizes---for example, $N = 106^3 \approx 1.2 \times 10^6$
points and $N_t = 85$ steps---the pre-computed projection reduces total
simulation time from $O(10^2\text{--}10^3)$~s to $O(10)$~s, an improvement of
one to two orders of magnitude.

\subsubsection{Conditions for Validity}

The pre-computed projection is exact whenever the body force is temporally
separable in the sense of Eq.~\eqref{eq:separable_force}. This condition is
satisfied in the standard ARF scenario where:
\begin{itemize}[leftmargin=1.5em]
\item the acoustic intensity field is computed once (steady-state or
      time-averaged over acoustic cycles),
\item the resulting radiation force spatial pattern $\bm{b}_0(\bm{x})$ is
      fixed, and
\item a scalar envelope $g(t)$ modulates the push on/off timing.
\end{itemize}
The approach is \emph{not} applicable when the spatial force distribution
itself changes with time, as would occur with a moving acoustic focus,
time-varying aberration correction, or forces derived from a non-stationary
pressure field.  In such cases, the general per-step CG projection of
Sec.~3.2 remains necessary.


\subsection{Boundary and initial conditions}
\label{sec:bc_ic}
In the current implementation, the computational domain is terminated with a fixed
(Dirichlet) boundary condition by setting displacement to zero on the outer layer.
This choice simplifies the staggered-grid update and provides a well-posed finite
domain problem, but it can introduce reflections that contaminate late-time arrivals.
Absorbing boundary treatments such as perfectly matched layers (PML)~\cite{berenger1994perfect,collino2001application} can reduce such artifacts when late-time fidelity is required. Accordingly, validation and analysis are typically restricted to early and
intermediate propagation times and/or to sensor offsets that are not dominated by
boundary interactions. The initial state is commonly specified as a divergence-free
velocity field constructed from a vector potential (a curl-based initialization),
which preferentially excites shear motion while suppressing compressional content.

\subsection{Validation approach and current implementation}
\label{sec:validation_methods}
Validation is treated as a layered process to separate implementation defects from
modeling limitations. At the most basic level, code-correctness checks should include
analytic or manufactured solutions in simplified geometries, symmetry checks (e.g.,
a symmetric source should produce a symmetric wavefield in a homogeneous medium), and
regression tests that detect unintended changes. Numerical-consistency checks then
assess stability and accuracy through grid and time-step refinement, sensitivity to
the Courant--Friedrichs--Lewy (CFL) condition~\cite{courant1928partial}, and characterization of numerical
dispersion and attenuation as resolution changes. Physics-consistency checks verify
that homogeneous-media propagation recovers the expected shear wavespeed
$c_s=\sqrt{\mu/\rho}$, that attenuation trends behave monotonically with the viscosity
parameter $\eta$ when present, and that passive cases do not exhibit non-physical
energy growth. Finally, boundary and geometry checks quantify sensitivity to domain
size and to the choice of boundary treatment (e.g., fixed versus absorbing layers).

The repository implements a focused wavespeed validation in
\texttt{fullwave/shear\_wave/validation/validate\_wavespeed.py}. A homogeneous benchmark medium is
defined by selecting a target shear speed $c_s$ and setting $\mu=\rho c_s^2$.
Shear-dominant excitation is obtained via curl-based initialization and a
divergence-free force projection. Displacement traces are recorded at multiple
offsets, and arrival times are estimated from a velocity proxy
$\partial u_z/\partial t$ using peak picking with distance-based start gating.
Recovered speed is reported per sensor, along with mean and linear-regression
estimates~\cite{urban2012review,chen2013quantitative}. This validation establishes correct recovery of shear-wave propagation
speed in controlled homogeneous cases~\cite{capuzzi2018numerical}, but does not by itself validate compressional
physics, anisotropy, strong heterogeneity, or broadband viscoelastic constitutive
behavior.

\section{Results}\label{sec:results}
\subsection{Non-Viscous Validation Run (\texorpdfstring{$\eta=0$}{eta=0}): Plots and Wavespeed Results}
To isolate elastic wavespeed recovery, the large-field validation was run with $\eta=0$ (no viscous term) on a cubic $180^3$ grid ($L_x=L_y=L_z=36$~mm, $\Delta x=0.2$~mm). The generated plots are shown below.

\begin{figure}[H]
\centering
\includegraphics[width=0.95\textwidth]{shear_slices_00320.png}
\caption{Representative shear-wave field snapshot ($u_z$, non-viscous run): XY slice (left), XZ slice (middle), YZ slice at $x=0$ (right). The YZ slice is near zero because the curl-based vector-potential initialization drives motion in the $x$-$z$ plane; the center $x$-plane is a nodal surface for $u_z$ by symmetry.}
\label{fig:shear_slices_xyz}
\end{figure}

\begin{figure}[H]
\centering
\includegraphics[width=0.85\textwidth]{shear_large_displacement.png}
\caption{Displacement traces ($u_z$) at increasing sensor offsets for the non-viscous validation run.}
\label{fig:nonvisc_disp}
\end{figure}

\begin{figure}[H]
\centering
\includegraphics[width=0.85\textwidth]{shear_large_velocity.png}
\caption{Velocity-proxy traces ($\partial u_z/\partial t$) used for arrival picking and wavespeed estimation.}
\label{fig:nonvisc_vel}
\end{figure}

Recovered values from the non-viscous run:
\begin{itemize}[leftmargin=1.5em]
\item Target speed: 2.000 m/s (from $\mu=\rho c_s^2$).
\item Sensor-wise estimates (m/s): 1.961, 1.976, 1.982, 1.986, 2.300, 2.603.
\item For validation metrics, the last two sensors are excluded (boundary/reflection sensitivity at long propagation distance).
\item Mean speed using first four sensors only: 1.976 m/s.
\item Linear fit speed using first four sensors only: 1.998 m/s.
\end{itemize}

\paragraph{Result interpretation.}
Using the first four sensors (nearest/intermediate range), the recovered speed is in close agreement with the imposed elastic target (\(2.000\) m/s): the mean is \(1.976\) m/s (\(\sim\)1.2\% low) and the distance-time linear-fit speed is \(1.998\) m/s (\(<0.1\%\) error). This supports that the shear-wave propagation physics and numerics are correct in the non-viscous benchmark regime. The two farthest sensors are excluded here because they are more sensitive to late-time boundary/reflection effects in the current finite domain; if needed, this can be reduced further by enlarging the domain or using absorbing boundaries.

\subsection{Grid Refinement / Convergence Study (Inviscid)}
An inviscid convergence sweep was performed with fixed physical domain size (\(L_x=L_y=L_z=36\) mm), fixed physical sensor distances (first four sensors only), and a six-level refinement range. Crucially, the Gaussian source width (\(\sigma_{xy}=5\) grid points) and velocity damping coefficient (\(\gamma_v=0.001\)) are held constant across all levels so that any differences are due purely to grid resolution.
\begin{itemize}[leftmargin=1.5em]
\item \(n=60\) (\(dX=6.00\times10^{-4}\) m): mean\(_4\)=1.828 m/s, \(c_{\mathrm{fit},4}=2.193\) m/s.
\item \(n=80\) (\(dX=4.50\times10^{-4}\) m): mean\(_4\)=1.907 m/s, \(c_{\mathrm{fit},4}=2.048\) m/s.
\item \(n=100\) (\(dX=3.60\times10^{-4}\) m): mean\(_4\)=1.938 m/s, \(c_{\mathrm{fit},4}=2.053\) m/s.
\item \(n=120\) (\(dX=3.00\times10^{-4}\) m): mean\(_4\)=1.954 m/s, \(c_{\mathrm{fit},4}=2.017\) m/s.
\item \(n=150\) (\(dX=2.40\times10^{-4}\) m): mean\(_4\)=1.968 m/s, \(c_{\mathrm{fit},4}=2.007\) m/s.
\item \(n=180\) (\(dX=2.00\times10^{-4}\) m): mean\(_4\)=1.972 m/s, \(c_{\mathrm{fit},4}=1.998\) m/s.
\end{itemize}

Two complementary error metrics are used to quantify convergence:
\begin{enumerate}
\item \textbf{Wavespeed error} (\(|\mathrm{mean}_4-c_{\mathrm{target}}|\)): the absolute difference between the mean recovered speed over the first four sensors and the analytical target. This is a scalar metric derived from discrete peak-picking, so it exhibits some variability.
\item \textbf{Trace L2 error}: the relative \(L_2\) norm of the velocity trace difference at each sensor, averaged over the four sensors, measured against the finest grid (\(n=180\)) as reference. This is a continuous metric that captures the full waveform convergence.
\end{enumerate}

\begin{figure}[H]
\centering
\includegraphics[width=0.78\textwidth]{convergence_mean4_error_only.png}
\caption{Wavespeed convergence: \( |\mathrm{mean}_4 - c_{\mathrm{target}}| \) vs grid spacing (log-log). Fitted slope 1.68, \(R^2=0.98\).}
\label{fig:convergence_mean4_only}
\end{figure}

\begin{figure}[H]
\centering
\includegraphics[width=0.78\textwidth]{convergence_trace_l2.png}
\caption{Trace L2 convergence: relative \(L_2\) error of velocity traces vs finest grid (log-log). Fitted slope 1.83, \(R^2=0.90\). Dashed reference lines show \(O(\Delta x)\) and \(O(\Delta x^2)\).}
\label{fig:convergence_trace_l2}
\end{figure}

\paragraph{Quantitative convergence metrics.}
Fitting
\(
|\mathrm{mean}_4-c_{\mathrm{target}}| \propto dX^p
\)
in log-log space gives a global slope of \(p=1.68\) (\(R^2=0.98\)), with pairwise observed orders:
\(
p_{60\to80}=2.17,\;
p_{80\to100}=1.77,\;
p_{100\to120}=1.64,\;
p_{120\to150}=1.63,\;
p_{150\to180}=0.76.
\)
The trace L2 metric gives a global slope of \(p=1.83\) (\(R^2=0.90\)), monotonically decreasing from \(1.83\) at \(n=60\) to \(0.32\) at \(n=150\).

\paragraph{Interpretation.}
Both metrics confirm convergence consistent with the second-order spatial discretization (7-point Laplacian stencil) and second-order leapfrog time integration. The trace L2 metric, which captures the full waveform rather than a single scalar, yields a cleaner slope (\(p\approx1.83\)) closer to the theoretical \(O(\Delta x^2)\) rate. The wavespeed metric (\(p\approx1.68\)) shows slightly sub-quadratic behavior due to the discrete nature of peak-picking for arrival time estimation. At the finest level (\(n=180\)), the linear-fit speed reaches \(1.998\) m/s (\(<0.1\%\) error).

\subsection{Viscosity Validation: 1D Harmonic Kelvin--Voigt Benchmark}
To avoid known bias from finite 3D pulse-domain effects, viscosity validation is reported here using only the dedicated 1D harmonic benchmark (\texttt{fullwave/shear\_wave/validation/validate\_kelvin\_voigt\_1d.py}), which better matches Kelvin--Voigt analytic assumptions (single frequency, effectively plane-wave propagation, homogeneous medium).

\subsubsection{1D harmonic benchmark results}
To align with Kelvin--Voigt theory and cover a brain-relevant range, the benchmark was run on \(\eta\in[0,1.5]\) Pa\,s with \(\Delta\eta=0.1\) Pa\,s. The high-\(\eta\) configuration uses (i) adaptive sensor placement based on the expected e-folding length, (ii) adaptive drive amplitude scaling, and (iii) a smooth frequency schedule \(f_0(\eta)\) guided by dimensional analysis.

Define the loss parameter and dimensionless attenuation
\begin{equation}
\mathrm{De}_v=\frac{\omega\eta}{\mu}, \qquad
\Pi_\alpha=\frac{\alpha c_s}{\omega}, \qquad c_s=\sqrt{\mu/\rho}.
\end{equation}
For weak-to-moderate loss (\(\mathrm{De}_v\ll 1\)),
\begin{equation}
\Pi_\alpha \approx \frac{\mathrm{De}_v}{2}
\;\Rightarrow\;
\alpha \approx \frac{\eta \omega^2}{2\rho c_s^3}.
\label{eq:alpha_scaling}
\end{equation}
Thus, to keep attenuation growth smooth with viscosity while staying tied to the Kelvin--Voigt asymptotic, we set a target power law
\begin{equation}
\alpha_{\mathrm{target}}(\eta)\propto \eta^p,
\end{equation}
which gives
\begin{equation}
f_0(\eta)\propto \eta^{-\beta}, \qquad \beta=\frac{1-p}{2}.
\end{equation}
In the implementation (\texttt{f0\_mode=dimensional}), \(p\) is set by \texttt{f0\_alpha\_exp} (default \(p=0.6\), hence \(\beta=0.2\)) unless overridden by \texttt{f0\_dim\_beta}. The provisional \(f_0(\eta)\) is then corrected by solving the exact Kelvin--Voigt dispersion relation so that exact \(\alpha(f_0,\eta)\) matches the asymptotic target. For \(\eta_{\max}=1.5\) Pa\,s with 800 Hz base drive and \texttt{f0\_min\_scale}=0.4, this yields a smooth drop to the 742 Hz range.

\paragraph{Adaptive sensor placement.}
A key challenge for attenuation measurement at high viscosity is that the e-folding length \(\ell_e=1/\alpha\) shrinks rapidly: at \(\eta=1.5\) Pa\,s the theoretical \(\alpha\approx 824\) Np/m gives \(\ell_e\approx 1.2\) mm. Sensors placed beyond a few e-folding lengths measure only noise, biasing the attenuation fit low.

To address this, sensor positions are adapted per \(\eta\) based on the theoretical attenuation at the scheduled drive frequency. For each case, \(\alpha_{\mathrm{theory}}(f_0,\eta)\) is computed from the exact Kelvin--Voigt dispersion relation, and \(n\) sensors are placed linearly from \(\tfrac{1}{2}\ell_e\) to \(3\ell_e\), ensuring all sensors lie within a measurable amplitude range. For the inviscid case (\(\eta=0\)), sensors fall back to a default spread covering 3\%--18\% of the domain.

\paragraph{Drive and frequency scheduling parameters.}
The run reported here used:
\begin{itemize}[leftmargin=1.5em]
\item Frequency scheduler: \texttt{f0\_mode=dimensional}, \(\eta_{\mathrm{ref}}=0.2\) Pa\,s, \(\texttt{f0\_alpha\_exp}=0.6\Rightarrow \beta=0.2\), smoothing factor \(0.7\), lower-bound scale \(\texttt{f0\_min\_scale}=0.4\).
\item Drive scaling: \(\eta_{\mathrm{hi}}=0.2\) Pa\,s, multiplier \(1+6(\eta-\eta_{\mathrm{hi}})\), capped at \(6\times\).
\item SNR diagnostics: \(\mathrm{SNR}_{\min}=6\) (used for sensor-quality flagging); attenuation figure shows the no-SNR-gate measured curve for visual continuity.
\end{itemize}

\begin{figure}[H]
\centering
\includegraphics[width=0.8\textwidth]{kelvin_voigt_1d_alpha.png}
\caption{1D harmonic Kelvin--Voigt benchmark: attenuation comparison (measured vs theory) over \(\eta\in[0,1.5]\) Pa\,s with adaptive sensor placement and frequency scheduling.}
\label{fig:kv1d_alpha}
\end{figure}

\begin{figure}[H]
\centering
\includegraphics[width=0.8\textwidth]{kelvin_voigt_1d_speed.png}
\caption{1D harmonic Kelvin--Voigt benchmark: phase-speed comparison (measured vs theory).}
\label{fig:kv1d_speed}
\end{figure}

Representative results from the 16-point viscosity sweep:
\begin{itemize}[leftmargin=1.5em]
\item \(\eta=0.10\) Pa\,s: \(\alpha_{\mathrm{meas}}=158.5\) Np/m vs \(\alpha_{\mathrm{th}}=156.4\) Np/m (\(+1.3\%\)); \(c_{\mathrm{meas}}=2.002\) m/s vs \(c_{\mathrm{th}}=2.012\) m/s (\(-0.5\%\)).
\item \(\eta=0.30\) Pa\,s: \(\alpha_{\mathrm{meas}}=436.8\) Np/m vs \(\alpha_{\mathrm{th}}=425.8\) Np/m (\(+2.6\%\)); \(c_{\mathrm{meas}}=2.083\) m/s vs \(c_{\mathrm{th}}=2.099\) m/s (\(-0.8\%\)).
\item \(\eta=0.50\) Pa\,s: \(\alpha_{\mathrm{meas}}=607.8\) Np/m vs \(\alpha_{\mathrm{th}}=589.5\) Np/m (\(+3.1\%\)); \(c_{\mathrm{meas}}=2.222\) m/s vs \(c_{\mathrm{th}}=2.239\) m/s (\(-0.8\%\)).
\item \(\eta=1.00\) Pa\,s: \(\alpha_{\mathrm{meas}}=809.7\) Np/m vs \(\alpha_{\mathrm{th}}=785.6\) Np/m (\(+3.1\%\)); \(c_{\mathrm{meas}}=2.729\) m/s vs \(c_{\mathrm{th}}=2.728\) m/s (\(<0.1\%\)).
\item \(\eta=1.50\) Pa\,s: \(\alpha_{\mathrm{meas}}=850.0\) Np/m vs \(\alpha_{\mathrm{th}}=824.2\) Np/m (\(+3.1\%\)); \(c_{\mathrm{meas}}=3.267\) m/s vs \(c_{\mathrm{th}}=3.281\) m/s (\(-0.4\%\)).
\item Monotonicity: both \(\alpha_{\mathrm{meas}}\) and \(\alpha_{\mathrm{th}}\) are monotone increasing over the full range \(\eta\in[0,1.5]\) Pa\,s.
\item Aggregate agreement over \(\eta\in[0.1,1.5]\) Pa\,s: mean absolute attenuation error \(\sim\)3\%, mean absolute phase-speed error \(<0.5\%\).
\end{itemize}

\paragraph{Conclusion.}
The combination of adaptive sensor placement (keyed to the theoretical e-folding length) and dimensional-analysis-guided frequency scheduling (Equation~\eqref{eq:alpha_scaling}) yields robust attenuation recovery across the full viscosity range. Both attenuation and phase speed track the analytical Kelvin--Voigt predictions with consistent \(\sim\)3\% and \(<0.5\%\) agreement, respectively, confirming the implicit viscous time-stepping scheme.



\section{Discussion}
\label{sec:discussion}

\subsection{Key physical assumptions}
The solver is derived under a small-strain linearization: displacement gradients are
assumed to be small and the stress--strain relationship is linear. This is typically
appropriate for ultrasound radiation-force excitations in SWE~\cite{sarvazyan1998elasticity,barr2015elastography}, where observed tissue
motion is commonly in the micrometer to tens-of-micrometers regime, but it can fail
under unusually strong pushes, large strains, near-contact compression, or when
tissue damage and nonlinear mechanics become relevant.

Material behavior is modeled as isotropic with an effectively homogeneous (or
piecewise homogeneous) shear response, represented by a scalar shear modulus
$\mu(\bm{x})$ and density $\rho(\bm{x})$~\cite{fung1993biomechanics,duck1990physical}. This approximation is often reasonable in
parenchymal soft tissues at the spatial resolution used for SWE processing~\cite{gennisson2013ultrasound}, but it
can break down in anisotropic or strongly heterogeneous structures such as skeletal
muscle, tendon, white-matter tracts, fibrotic septa, layered media, and near rigid
interfaces (e.g., skull or other hard boundaries).

A central modeling choice is to represent shear-dominant propagation while omitting
explicit compressional (longitudinal) elastodynamics. This is motivated by the
near-incompressibility of soft tissue ($K\gg\mu$), which yields a strong separation
between longitudinal and shear wave speeds ($c_L\gg c_S$) and hence their relevant
timescales. In many SWE workflows, fast compressional transients exit the analysis
window quickly or are reduced by directional and temporal filtering, whereas the
quantity of interest is the slower transverse shear wave used for wavespeed and
modulus estimation. The tradeoff is that this reduced model does not capture
near-source pressure-wave content, P--S mode conversion at interfaces or strong
heterogeneities, nonlinear shear or shear shock dynamics~\cite{tripathi2017shearshock,espindola2017shearshock,tripathi2019shearshock2d}, or any diagnostic relying on volumetric strain.

Viscous attenuation is represented using a Kelvin--Voigt-like shear viscosity~\cite{catheline2004viscoelastic,ferry1980viscoelasticity}, in
which damping enters through a Laplacian term $\eta\nabla^2\dot{\bm{u}}$. This
single-parameter form is useful for capturing qualitative attenuation and pulse
broadening, but soft tissue often exhibits frequency-dependent, power-law, or
fractional viscoelastic behavior~\cite{urban2012review}; consequently, a constant $\eta$ may be
insufficient for broadband inversions or dispersion-sensitive analyses.

Finally, the present implementation enforces fixed (zero-displacement) boundaries.
This simplifies numerics but can generate reflections that bias late-time arrival
picking. Practical mitigation strategies include using larger computational domains,
restricting analysis to earlier times/shorter propagation distances, or implementing
absorbing layers (e.g., PML)~\cite{berenger1994perfect,collino2001application} when late-time fidelity is required.



\subsection{Connection to ultrasound-driven shear rheology and SWE}
In ARF-based SWE~\cite{nightingale2002radiation,bercoff2004supersonic}, the acoustic field produces a predominantly axial body force in the
push region. The resulting mechanical response contains both (i) a rapid local
compression-related transient near the source and (ii) slower transverse shear-wave
radiation away from the push region. When the scientific or clinical objective is
shear-wave arrival timing and wavespeed in soft tissue, the reduced governing
equation used here provides an efficient approximation that isolates the dominant
shear dynamics. Conversely, when early-time fields, compressional contamination,
mode conversion, poroelastic fluid--solid coupling, or complex boundary interactions
are central to the question, a full elastodynamic or poroelastic model is required.

\subsection{When the assumptions are likely valid, and when they are likely invalid}
The reduced model is most appropriate in regimes with small displacements and
approximately linear response~\cite{sarvazyan1998elasticity}; in nearly incompressible soft tissue where SWE
biomarkers depend primarily on the shear-wave speed~\cite{sandrin2003transient,tanter2008quantitative,barr2015elastography}; in moderate propagation ranges
where boundary reflections are not dominant; and in analyses that are effectively
narrowband such that a Kelvin--Voigt damping term provides an adequate attenuation
approximation~\cite{catheline2004viscoelastic}.

The assumptions are more likely to fail in strongly anisotropic or heterogeneous
media (e.g., muscle fibers, scar bands, skull-coupled configurations), in problems
where compressional waves must be represented explicitly (especially near-source
early-time dynamics), in regimes where nonlinear shear or shear shock formation
is significant~\cite{tripathi2017shearshock,espindola2017shearshock}, in broadband rheological inversions requiring more faithful
frequency-dependent constitutive models, in domains where fixed-wall reflections
strongly influence arrival times, and in settings where poroelasticity or other
multiphase physics is needed to represent fluid--solid relative motion.

\subsection{Practical interpretation of inputs and outputs}
The solver outputs displacement and velocity fields that should be interpreted as
shear-dominant wavefields. In homogeneous media, the characteristic shear-wave speed
is $c_S\approx\sqrt{\mu/\rho}$, with practical bias introduced by viscosity,
heterogeneity, finite-domain effects, and the specific arrival-picking strategy.
In ultrasound multiphysics workflows, the body-force input should be understood as
a radiation-force–equivalent source term~\cite{palmeri2011radiation,rooney1973radiation} derived from an acoustic simulation~\cite{jensen1996field,treeby2010modeling,pinton2020fullwave} (or
from a prescribed analytic approximation), applied on the elastic timescale after
the acoustic field has effectively traversed the domain. This separation of scales
supports tight integration with acoustic tools while maintaining computational
efficiency for shear dynamics.

\subsection{Typical acoustic radiation force magnitudes}
\label{sec:arf_magnitudes}

The body force that drives the shear solver is the acoustic radiation force
(ARF), given by the plane-wave approximation~\cite{nyborg1965acoustic,torr1984acoustic}
\begin{equation}
  b = \frac{2\,\alpha_{\mathrm{Np}}}{c}\,I_{\mathrm{SPPA}},
  \qquad
  I_{\mathrm{SPPA}} = \frac{\langle p^2\rangle}{\rho\,c},
  \label{eq:arf_plane}
\end{equation}
where $\alpha_{\mathrm{Np}}$ is the absorption in Np/m at the transmit
frequency $f_0$, and $I_{\mathrm{SPPA}}$ is the spatial-peak
pulse-average intensity at the focus.  Because absorption in soft tissue
scales approximately linearly with
frequency~\cite{szabo2014diagnostic,duck1990physical}, the ARF magnitude
depends strongly on $f_0$ even at the same focal intensity.
Table~\ref{tab:arf_scaling} illustrates this effect for a representative
tissue absorption of 0.5~dB\,cm$^{-1}$\,MHz$^{-1}$ and a focal
$I_{\mathrm{SPPA}}$ of 1000~W/cm$^2$.

\begin{table}[H]
\centering
\caption{Frequency dependence of ARF body force at fixed focal intensity
  ($I_{\mathrm{SPPA}} = 1000$~W/cm$^2$,
   $\alpha_0 = 0.5$~dB\,cm$^{-1}$\,MHz$^{-1}$, $c = 1540$~m/s).}
\label{tab:arf_scaling}
\begin{tabular}{ccc}
\hline
$f_0$ (MHz) & $\alpha_{\mathrm{Np}}$ (Np/m) & $b$ (N/m$^3$) \\
\hline
1   &  5.8  & 750   \\
3   & 17.3  & 2\,250  \\
5   & 28.8  & 37\,400 \\
6.7 & 38.6  & 50\,100 \\
\hline
\end{tabular}
\end{table}

\noindent
In clinical soft-tissue ARFI at 5--7~MHz, focal intensities of
500--2000~W/cm$^2$ and correspondingly high attenuation produce body
forces on the order of $10^4$--$10^5$~N/m$^3$, yielding tissue
displacements of
1--20~$\mu$m~\cite{nightingale2002radiation,palmeri2005fem,palmeri2011radiation}.

For \emph{transcranial} applications at lower frequencies (0.5--1~MHz), the
situation differs in two important ways.  First, the absorption per unit
path length is roughly an order of magnitude lower than at diagnostic
frequencies, so the ARF per unit intensity is proportionally weaker.
Second, the skull attenuates and aberrates the transmitted beam, reducing
the focal intensity to roughly 30--50\% of its free-field
value~\cite{pinton2012braintherapy}.  The net effect is that body forces at
1~MHz are typically $10^2$--$10^4$~N/m$^3$ (two orders of magnitude below
high-frequency SWE), and the resulting brain tissue displacements are
sub-micrometer to a few micrometers.  This is consistent with transcranial
MR-ARFI measurements reporting 0.5--1.2~$\mu$m displacements in primate
brain at comparable acoustic parameters.

The transcranial demonstration in Appendix~\ref{app:transcranial_demo}
provides a concrete example: a 1024-element sparse array at 1~MHz with
1.5~MPa source pressure produces a peak focal pressure of $\sim$5~MPa,
a focal $I_{\mathrm{SPPA}}$ of $\sim$633~W/cm$^2$ (after CW envelope
correction), and a peak ARF of
$\sim$41{,}800~N/m$^3$ after skull propagation, yielding $\sim$1.7~$\mu$m
displacement in brain tissue ($\mu \approx 4.6$~kPa, $\eta = 0.5$~Pa$\cdot$s).

\subsection{From tissue displacement to membrane tension}
In neuromodulation applications, the quantity that ultimately gates
mechanosensitive ion channels is not tissue displacement itself but the
\emph{membrane tension} $T$ (force per unit length, mN/m) experienced by
the lipid bilayer in which the channel is embedded.  Bridging the
macroscopic displacement field produced by this solver to membrane tension
requires two steps: (i)~computing the tissue-level strain tensor from the
displacement gradient, and (ii)~converting the deviatoric strain to an
effective bilayer tension via the membrane area-expansion modulus.

The linearized strain tensor $\varepsilon_{ij} = \tfrac{1}{2}(\partial_i u_j + \partial_j u_i)$
is computed from the solver displacement field $\bm{u}(\bm{x},t)$.  Because
the tissue is nearly incompressible, the volumetric (dilatational) strain is
small and the mechanically relevant component is the deviatoric strain norm
$\varepsilon_{\mathrm{dev}}$.  The effective membrane tension is then
\begin{equation}
    T = K_A \,\alpha\, \varepsilon_{\mathrm{dev}} \,,
    \label{eq:tension}
\end{equation}
where $K_A$ is the area-expansion modulus of the lipid bilayer
(typically 0.1--0.5~N/m for biological
membranes~\cite{rawicz2000effect}) and $\alpha$ is a dimensionless
coupling factor relating far-field tissue strain to local membrane area
change.

The tension $T$ enters the Boltzmann gating model for mechanosensitive
channels.  Each channel opens with a steady-state probability
$P_o(T) = [1 + \exp(-\Delta A\,(T - T_{1/2})/k_BT_{\mathrm{abs}})]^{-1}$,
where $\Delta A$ is the effective gate area change and $T_{1/2}$ is the
half-activation tension---the tension at which the channel is open 50\%
of the time.

$T_{1/2}$ is determined experimentally.  Two methods are used in practice.
The most common is \emph{patch-clamp electrophysiology with applied suction}:
a glass micropipette is sealed onto a membrane patch, negative pressure
(mmHg) is applied, and channel currents are recorded as a function of
pressure.  Converting the applied pressure $P$ to tension requires the
Laplace relation $T = Pr/2$, where $r$ is the radius of curvature of the
membrane dome inside the pipette.  Because $r$ depends on pipette geometry
and seal quality and is rarely measured directly, this conversion introduces
uncertainty of order $\pm$1--2~mN/m~\cite{lewis2015mechanical}.  The more
direct approach uses \emph{reconstituted lipid bilayers}, in which purified
channel proteins are inserted into artificial membranes whose tension is
controlled and measured directly, eliminating the Laplace
ambiguity~\cite{syeda2016piezo,brohawn2014mechanosensitivity}.

For the channels most relevant to ultrasonic neuromodulation, the best-characterized
half-activation tensions are: Piezo1, $T_{1/2} = 2.7 \pm 0.4$~mN/m, measured by
direct tension clamp with simultaneous patch imaging~\cite{lewis2015mechanical}
and confirmed in cell-free bilayers~\cite{syeda2016piezo}; and TRAAK,
$T_{1/2} \approx 5$~mN/m, from reconstituted
bilayers~\cite{brohawn2014mechanosensitivity}.  Values for TREK-1
($\sim$3.5~mN/m) and TREK-2 ($\sim$4~mN/m) carry larger uncertainty because
they are derived from patch-clamp pressure data with assumed pipette
geometry~\cite{patel1998lipid,bang2000trek2}.

These tension scales set the target for the simulation pipeline: the shear
wave solver must resolve displacement gradients (and hence strains)
accurately enough that the resulting membrane tensions can be compared with
the channel $T_{1/2}$ values.  For a representative $K_A = 0.25$~N/m and
$\alpha = 1$, a deviatoric strain of $\sim$10$^{-5}$ produces a membrane
tension of $\sim$2.5~$\mu$N/m---well below channel thresholds---while strains
of order $10^{-2}$ reach the few-mN/m range where channel gating becomes
significant.  This underscores the importance of accurate strain
computation from the displacement field and motivates the use of the
deviatoric invariant (which isolates shape-changing deformation) rather
than the full strain tensor.


\section{Conclusion}
\label{sec:conclusion}
This manuscript derived the reduced shear elastodynamic equations implemented in
\texttt{fullwave/shear\_wave/shear\_fdtd\_staggered\_3d.py} and documented the physical assumptions under
which the formulation is appropriate for ultrasound-driven mechanics, including SWE~\cite{gennisson2013ultrasound,sigrist2017ultrasound,palmeri2011radiation}
and related elastomechanical hypotheses in neuromodulation~\cite{blackmore2019ultrasound,jones2022neuromodulation,pinton2012braintherapy}. The resulting solver
targets shear-dominant motion in nearly incompressible soft solids, enabling
computationally efficient simulation of radiation-force–induced displacements while
remaining compatible with sequential acoustic--elastic coupling strategies that
reflect the strong separation of longitudinal and shear timescales.

Validation in homogeneous benchmark media demonstrates accurate recovery of the
imposed shear wavespeed under controlled conditions, and convergence sweeps support
numerical consistency over practical grid resolutions. The principal limitations are
those implied by the model reduction and boundary treatment: compressional dynamics,
mode conversion at interfaces, anisotropy, strong heterogeneity, broadband
frequency-dependent viscoelasticity, and late-time reflection artifacts are not
fully represented. Within its intended scope, the reduced solver provides a
lightweight, purpose-built tool for multiphysics ultrasound workflows where shear
wave generation, propagation, and measurement are the dominant quantities of
interest~\cite{sarvazyan1998elasticity,bercoff2004supersonic,barr2015elastography,nightingale2002radiation,tanter2008quantitative}.

\bibliographystyle{plain}
\bibliography{general}

\appendix

\section{End-to-End Transcranial Shear Wave Demo}
\label{app:transcranial_demo}

This appendix presents an end-to-end demonstration of the complete
acoustic--elastic pipeline applied to transcranial focused ultrasound through a
realistic human skull. The simulation uses a 1024-element sparse phased array
focused through an \emph{ex~vivo} micro-CT skull (Halle temporal bone specimen) to
generate acoustic radiation force (ARF) at a 50~mm focal depth, which then drives
shear wave displacement in brain tissue via the FDTD solver derived in this
document.

The pipeline proceeds in four stages:
\begin{enumerate}[nosep]
  \item Sparse array acoustic transmit through micro-CT skull (GPU FDTD),
  \item Computation of the acoustic radiation force from the pressure field,
  \item Shear wave FDTD simulation of viscoelastic brain tissue displacement,
  \item Visualization of all intermediate and final fields.
\end{enumerate}

\noindent
The corresponding script is
\texttt{examples/wave\_3d/sparse\_transcranial\_shear\_3d.py}.

\subsection{Sparse Array Geometry}

The transducer is a 1024-element sparse phased array with 750~$\mu$m element pitch.
Element positions are loaded from an HDF5 connector file and mapped onto the
simulation grid ($66 \times 66$~mm lateral--elevational aperture). Per-element
focusing delays steer the beam to a 50~mm depth focal point.
Figure~\ref{fig:array_surface} shows the element layout in the transducer plane.

\begin{figure}[H]
  \centering
  \includegraphics[width=0.6\textwidth]{fig_a_array_surface.png}
  \caption{Sparse array surface layout showing element center positions in the
    lateral--elevational ($y$--$z$) plane.}
  \label{fig:array_surface}
\end{figure}

\subsection{Acoustical Medium Properties}

The heterogeneous skull medium is constructed from a micro-CT NRRD volume of the
Halle temporal bone specimen using the \texttt{MicroCTSkullDomain} preset. Hounsfield
units are mapped to spatially varying sound speed (1540--2900~m/s), density
(1000--2200~kg/m$^3$), and attenuation coefficient. The background tissue is
assigned standard soft-tissue properties. Figure~\ref{fig:acoustical_maps} shows
the resulting medium property maps.

\begin{figure}[H]
  \centering
  \includegraphics[width=\textwidth]{fig_b_acoustical_maps.png}
  \caption{Acoustical medium properties derived from micro-CT skull data.
    \textbf{Top row:} $x$--$y$ (depth--lateral) slices at center elevation showing
    sound speed, density, and attenuation coefficient $\alpha_0$.
    \textbf{Bottom row:} $x$--$z$ (depth--elevational) slices at center lateral
    position. The skull layer is clearly visible as a high-speed, high-density
    region near the transducer surface.}
  \label{fig:acoustical_maps}
\end{figure}

\subsection{Acoustic Pressure Propagation}

The acoustic simulation uses the Fullwave GPU FDTD solver with exponential
attenuation at $f_0 = 1$~MHz, source pressure amplitude of 1.5~MPa (typical
clinical ARFI), and a sparse-grid sensor (4$\times$ downsampling in each spatial
dimension). Figure~\ref{fig:pressure_snapshots} shows the focused beam propagating
through the skull into brain tissue.

\begin{figure}[H]
  \centering
  \includegraphics[width=\textwidth]{fig_c_pressure_snapshots.png}
  \caption{Pressure propagation snapshots in the $x$--$y$ plane at center
    elevation. The focused beam emerges from the transducer (left edge),
    propagates through the skull (black contour), converges at the 50~mm focal
    zone, and subsequently diverges. Five time frames are shown spanning the
    full propagation duration.}
  \label{fig:pressure_snapshots}
\end{figure}

\subsection{Acoustic Radiation Force}

Two formulations for the acoustic radiation force (ARF) are available.

\paragraph{Pressure-based (plane-wave) approximation.}
The simplest approach assumes locally planar wavefronts and computes a scalar
body force along the beam axis:
\begin{equation}
  b_0 = \frac{2\,\alpha_{\mathrm{Np}}}{c}\,I_{\mathrm{SPPA}},
  \qquad
  I_{\mathrm{SPPA}} = \frac{1}{\rho\, c}
    \;\frac{\sum_{i=1}^{N_{\mathrm{total}}} p_i^2}{N_{\mathrm{pulse}}},
  \label{eq:isppa_norm}
\end{equation}
where the sum runs over all $N_{\mathrm{total}}$ recorded time steps of the
transient simulation and is normalized by the number of source-active steps
$N_{\mathrm{pulse}}$.  This \emph{silence correction} accounts for the dead
time after the pulse has passed through the domain.  A second
\emph{CW~envelope correction} is needed because the Gaussian-windowed
source pulse has a lower time-averaged $\langle p^2 \rangle$ than a
continuous-wave signal at the same peak pressure.  For CW transmission,
$\langle p^2 \rangle = p_{\mathrm{peak}}^2/2$, whereas for a short
Gaussian-modulated pulse the mean is much smaller.  The correction factor
\begin{equation}
  C_{\mathrm{CW}} = \frac{p_{\mathrm{peak}}^2 / 2}
                          {\langle p_{\mathrm{src}}^2 \rangle}
  \label{eq:cw_correction}
\end{equation}
scales the pulse-average intensity up to the equivalent CW value.  For a
3-cycle pulse ($C_{\mathrm{CW}} \approx 4.1$) and a 2-cycle pulse
($C_{\mathrm{CW}} \approx 7.2$).  The total normalization factor is the
product of the silence and CW corrections.
The same normalization is used for the volumetric heat source
$Q = 2\,\alpha_{\mathrm{Np}}\,I_{\mathrm{SPPA}}$ in the bioheat module,
ensuring a consistent mapping from transient simulations to equivalent
continuous-wave intensity fields.

\paragraph{Velocity-based (Poynting vector) formulation.}
For converging beams the plane-wave assumption breaks down near the focal zone,
where wavefronts are curved and energy flux has significant lateral components.
A more accurate approach uses the time-averaged Poynting vector
$\bm{I} = \langle p\, \bm{v} \rangle$, where $\bm{v}$ is the particle
velocity field obtained directly from the FDTD solver
(\texttt{output\_velocity=True}).  The vector body force is then
\begin{equation}
  f_i = \frac{2\,\alpha_{\mathrm{Np}}}{c}\,I_i,
  \qquad
  I_i = \langle p\, v_i \rangle,
  \qquad i \in \{x,y,z\},
  \label{eq:arf_poynting}
\end{equation}
yielding three spatially varying force components $(f_x, f_y, f_z)$.

In this demonstration we use the velocity-based formulation
(\texttt{-{}-use-velocity}) with the combined silence and CW envelope
corrections (Eqs.~\ref{eq:isppa_norm}--\ref{eq:cw_correction}).
The radiation force within the skull is set to zero in both formulations (bone is
rigid at the shear wave scale).
Figure~\ref{fig:radiation_force} shows the depth-component ($f_x$) of the
resulting body force distribution.

\paragraph{Results.}
With the combined pulse normalization (silence $\times$ CW correction,
total factor $\approx 57.8\times$ for this 2-cycle pulse), the
velocity-based depth component gives
$f_x^{\max} \approx 41{,}800$~N/m$^3$, with lateral and elevational
components $|f_y|^{\max} \approx 11{,}800$~N/m$^3$ and
$|f_z|^{\max} \approx 9{,}750$~N/m$^3$.  The off-axis components represent
roughly 25\% of the peak axial force and are passed to the shear wave
solver as three-component driving terms.

\begin{figure}[H]
  \centering
  \includegraphics[width=\textwidth]{fig_d_radiation_force.png}
  \caption{Acoustic radiation force (ARF) body force density (depth component
    $f_x$, velocity-based Poynting vector formulation with pulse-length and
    CW envelope normalization).
    \textbf{Left:} $x$--$y$ slice. \textbf{Right:} $x$--$z$ slice. The force is
    concentrated at the focal spot ($\sim$41{,}800~N/m$^3$ peak), with the skull
    region masked to zero.}
  \label{fig:radiation_force}
\end{figure}

\subsection{Shear Wave Displacement}

The shear wave simulation uses the staggered-grid FDTD solver (Section~3) with
brain tissue parameters derived from the porcine brain measurements of
Esp\'{i}ndola et al.~\cite{espindola2017shearshock} and Chandrasekaran
et al.~\cite{chandrasekaran2021modeling}: shear wave speed
$c_s = 2.14$~m/s at 75~Hz, density $\rho = 1000$~kg/m$^3$, giving shear
modulus $\mu = \rho\,c_s^2 \approx 4580$~Pa and Young's modulus
$E = 3\mu \approx 13.7$~kPa (assuming Poisson's ratio $\nu = 0.5$).
Kelvin--Voigt viscosity is set to $\eta = 0.5$~Pa$\cdot$s. All three ARF
components $(f_x, f_y, f_z)$ from the velocity-based formulation are applied
as forcing terms via the separable temporal envelope (\texttt{env\_t}), enabling
a single pre-computed Helmholtz projection.

Figure~\ref{fig:displacement_snapshots} shows the shear wave displacement
magnitude $|\mathbf{u}|$ at five time points, demonstrating the wave expanding
radially from the focal zone.

\begin{figure}[H]
  \centering
  \includegraphics[width=\textwidth]{fig_e_displacement_snapshots.png}
  \caption{Shear wave displacement magnitude snapshots.
    \textbf{Top row:} $x$--$y$ slices at center elevation.
    \textbf{Bottom row:} $x$--$z$ slices at center lateral position.
    The displacement field expands radially from the focal spot over
    0.7--3.0~ms. White contours indicate the skull boundary.}
  \label{fig:displacement_snapshots}
\end{figure}

Figure~\ref{fig:displacement_3d} provides a 3D rendering of the displacement
magnitude, showing the spatial extent of the shear wave field.

\begin{figure}[H]
  \centering
  \includegraphics[width=0.7\textwidth]{fig_f_displacement_3d.png}
  \caption{3D rendering of the displacement magnitude field. The isosurface
    at 50\% of peak $|\mathbf{u}|$ is shown in coral, with the skull rendered
    as a translucent gray surface. The focal displacement region is clearly
    localized in the brain tissue volume behind the skull.}
  \label{fig:displacement_3d}
\end{figure}

Figure~\ref{fig:displacement_trace} shows the temporal evolution of the
displacement at the peak radiation force location, illustrating the pulsed
tissue response to a single 2~ms push followed by 8~ms of viscous relaxation
(20\% duty cycle, 10~ms total).  Peak displacement reaches
$|u_x|^{\max} \approx 1.7$~$\mu$m.

\begin{figure}[H]
  \centering
  \includegraphics[width=0.85\textwidth]{fig_g_displacement_trace.png}
  \caption{Displacement $u_x$ at the peak ARF location as a function of time.
    The gray shading indicates the force envelope (single 2~ms push).
    The tissue displacement rises during the forcing interval and relaxes
    viscously during the subsequent 8~ms off period.}
  \label{fig:displacement_trace}
\end{figure}

\subsection{Theta-Burst TUS Protocol}
\label{sec:theta_burst}

The pipeline also supports a theta-burst transcranial ultrasound stimulation
(TUS) protocol motivated by neuromodulation
applications~\cite{deffieux2013ultrasoundneuromod,legon2014transcranial}.
Theta-burst TUS adapts the theta-burst pattern used in transcranial magnetic
stimulation (TMS) to focused ultrasound, delivering brief acoustic pulses
grouped at the 5~Hz theta rhythm that is prominent in hippocampal and cortical
circuits.

\paragraph{Protocol parameters.}
The theta-burst envelope consists of 360~$\mu$s tone-bursts repeated at a
1~kHz pulse repetition frequency (PRF), grouped into trains of 3~bursts.
Each 3-burst group occupies the first 3~ms of a 200~ms theta period
($1/5$~Hz), giving an overall duty cycle of $\sim$0.5\%.  One complete
theta cycle (200~ms) is simulated.

\paragraph{Intensity scaling.}
The acoustic simulation described in Section~\ref{sec:theta_burst} was
performed at a source pressure of 1.5~MPa, yielding a measured focal
spatial-peak pulse-average intensity of $I_{\mathrm{SPPA}} \approx 633$~W/cm$^2$
(after CW envelope correction).
Typical neuromodulation protocols operate at lower intensities,
$I_{\mathrm{SPPA}} \approx 20$--$30$~W/cm$^2$~\cite{deffieux2013ultrasoundneuromod,legon2014transcranial}.
Rather than re-running the acoustic simulation, the body force is scaled by
the ratio of target-to-simulated focal intensity:
\begin{equation}
  \bm{f}_{\mathrm{neuro}} = \frac{I_{\mathrm{SPPA}}^{\mathrm{target}}}
                                  {I_{\mathrm{SPPA}}^{\mathrm{sim}}}
                            \;\bm{f}_{\mathrm{sim}},
  \label{eq:arf_scaling}
\end{equation}
since the radiation force is proportional to acoustic intensity.  For a
target of 30~W/cm$^2$ the scaling factor is $30/633 \approx 0.047$, reducing
the peak axial force from $\sim$41{,}800 to $\sim$2{,}000~N/m$^3$.

\paragraph{Results.}
With the scaled ARF and theta-burst envelope, the shear FDTD solver
(12\,035 time steps over 200~ms) produces a peak displacement of
$|u_x|^{\max} \approx 0.3$~nm at the end of the theta cycle.
Instantaneous displacement during the 3~ms burst window reaches
$\sim$0.08~$\mu$m before viscous damping ($\eta = 0.5$~Pa$\cdot$s)
dissipates the motion during the subsequent 197~ms off-period.  The
sub-nanometer residual at cycle end is consistent with the very low duty
cycle (0.5\%) and confirms that each theta cycle effectively returns the
tissue to rest before the next group arrives.

\paragraph{Physical interpretation.}
The displacement magnitudes are orders of magnitude below those produced by
clinical ARFI pushes (typically 1--20~$\mu$m), which is expected:
neuromodulation operates at much lower intensities and duty cycles.
The neuromodulatory mechanism is not believed to depend on large tissue
displacements but rather on the mechanical perturbation of ion channels
and membrane dynamics at the cellular
scale~\cite{blackmore2019ultrasound,jones2022neuromodulation}.  The
theta-burst temporal structure provides a means to entrain neural
oscillatory activity at the theta frequency while keeping time-averaged
intensity low for safety.

\subsection{Summary}

This demonstration validates the complete sequential pipeline---acoustic transmit,
radiation force computation, and shear wave FDTD---in a realistic transcranial
geometry.  The combined pulse-length and CW envelope normalization
(Eqs.~\ref{eq:isppa_norm}--\ref{eq:cw_correction}) provides a
consistent mapping from the transient acoustic simulation to equivalent
continuous-wave intensity, shared by both the radiation force and bioheat
modules.  The velocity-based Poynting vector ARF formulation
(Eq.~\ref{eq:arf_poynting}) captures off-axis energy flux near the focal zone,
yielding significant lateral and elevational force components ($\sim$25\% of the
peak axial force).  A single 2~ms push produces a peak displacement of
$\sim$1.7~$\mu$m, with the shear wave expanding radially from the focal zone
at the theoretical shear wave speed.

The theta-burst TUS mode (Section~\ref{sec:theta_burst}) extends the pipeline
to neuromodulation applications by scaling the body force to a target focal
$I_{\mathrm{SPPA}}$ and applying a physiologically motivated temporal envelope
(360~$\mu$s bursts at 1~kHz PRF, grouped at 5~Hz theta).

The pipeline is implemented in a single script with
command-line options for skipping previously computed stages
(\texttt{-{}-skip-acoustic}, \texttt{-{}-skip-shear}), selecting the ARF
formulation (\texttt{-{}-use-velocity}), and enabling the theta-burst
protocol (\texttt{-{}-theta-burst}), facilitating iterative parameter
studies.

\section{TIPS Transcranial Focused Ultrasound Demo}
\label{app:tips_demo}

This appendix presents a second end-to-end demonstration using the Philips
TIPS (Therapeutic Imaging and Phased~array Sonication) transducer, a
curved-bowl annular phased array designed for transcranial focused
ultrasound therapy and neuromodulation.  The TIPS geometry differs
fundamentally from the sparse phased array of
Appendix~\ref{app:transcranial_demo}: rather than discrete point-like
elements on a flat aperture, the TIPS uses a spherical-cap bowl surface
with concentric annular ring elements, producing a geometrically focused
beam that is then electronically steered via per-element delays.

The pipeline follows the same four-stage sequential workflow:
\begin{enumerate}[nosep]
  \item TIPS curved-bowl acoustic transmit through micro-CT skull (GPU FDTD),
  \item Computation of acoustic radiation force from the pressure field,
  \item Shear wave FDTD simulation of viscoelastic brain tissue displacement,
  \item Visualization of all intermediate and final fields.
\end{enumerate}

\noindent
The corresponding script is
\texttt{examples/wave\_3d/tips\_transcranial\_shear\_3d.py}.

\subsection{TIPS Transducer Geometry}

The TIPS transducer is modeled as a spherical-cap bowl with the following
parameters: radius of curvature $R = 80$~mm, inner aperture radius
$r_{\mathrm{in}} = 20.5$~mm (hollow center), outer aperture radius
$r_{\mathrm{out}} = 46$~mm (92~mm full diameter), 8~concentric annular
elements of equal radial width, and a bowl thickness of 3~voxel layers
along the surface normal.  Per-pixel geometric focusing delays steer the
beam to a 50~mm depth focal point, with delays averaged within each
annular element to approximate the hardware phasing of the physical
transducer.

Figure~\ref{fig:tips_source_mask} shows the TIPS bowl source mask in the
central elevational slice, illustrating the concentric annular geometry.

\begin{figure}[H]
  \centering
  \includegraphics[width=0.6\textwidth]{fig_a_tips_source_mask.png}
  \caption{TIPS bowl source mask in the $x$--$y$ plane (central
    elevational slice).  The concentric annular ring structure and hollow
    center are clearly visible.  The bowl curvature produces a crescent
    cross-section in the depth--lateral plane.}
  \label{fig:tips_source_mask}
\end{figure}

\subsection{Acoustical Medium and Pressure Propagation}

The heterogeneous skull medium uses the same Halle temporal bone micro-CT
specimen and \texttt{MicroCTSkullDomain} preset as
Appendix~\ref{app:transcranial_demo}.  Because the TIPS bowl surface
extends further in depth than a flat aperture, the skull slab is shifted
deeper to avoid geometric intersection between the transducer source
voxels and the bone.

The acoustic simulation uses a source pressure amplitude of 400~kPa
(lower than the 1.5~MPa used for the sparse array, reflecting
neuromodulation operating levels) with 3-cycle Gaussian-modulated
sinusoidal pulses at $f_0 = 1$~MHz.  The domain size is
$80 \times 112 \times 112$~mm with 10~points per wavelength.

Figure~\ref{fig:tips_acoustical_maps} shows the acoustical medium property
maps, and Figure~\ref{fig:tips_pressure_snapshots} shows the focused beam
propagating through the skull.

\begin{figure}[H]
  \centering
  \includegraphics[width=\textwidth]{fig_tips_b_acoustical_maps.png}
  \caption{Acoustical medium properties for the TIPS transcranial
    configuration.  \textbf{Top row:} $x$--$y$ slices at center elevation.
    \textbf{Bottom row:} $x$--$z$ slices at center lateral position.  The
    skull layer appears as the high-speed, high-density region, shifted
    deeper relative to the sparse array configuration to clear the TIPS
    bowl geometry.}
  \label{fig:tips_acoustical_maps}
\end{figure}

\begin{figure}[H]
  \centering
  \includegraphics[width=\textwidth]{fig_tips_c_pressure_snapshots.png}
  \caption{Pressure propagation snapshots for the TIPS focused transmit.
    The geometrically focused bowl produces a converging wavefront that
    propagates through the skull and focuses at 50~mm depth.  Compared to
    the sparse array (Figure~\ref{fig:pressure_snapshots}), the TIPS bowl
    produces a smoother wavefront due to its continuous annular aperture.}
  \label{fig:tips_pressure_snapshots}
\end{figure}

\subsection{Acoustic Radiation Force}

The acoustic radiation force is computed using the plane-wave
approximation with the same pulse-length and CW envelope normalization
(Eqs.~\ref{eq:isppa_norm}--\ref{eq:cw_correction}) as the sparse array demo.
Radiation force within the skull is masked to zero.
The peak body force after skull masking is
$b_0^{\max} \approx 22{,}600$~N/m$^3$.
Figure~\ref{fig:tips_radiation_force} shows the resulting body force
distribution.

\begin{figure}[H]
  \centering
  \includegraphics[width=\textwidth]{fig_tips_d_radiation_force.png}
  \caption{Acoustic radiation force (depth component $b_0$) for the TIPS
    focused transmit.  \textbf{Left:} $x$--$y$ slice.  \textbf{Right:}
    $x$--$z$ slice.  The force is concentrated at the 50~mm focal spot,
    with the skull region masked to zero.}
  \label{fig:tips_radiation_force}
\end{figure}

\subsection{Shear Wave Displacement}

The shear wave simulation uses the same brain tissue parameters as the
sparse array demo (Appendix~\ref{app:transcranial_demo}).
The pulsed forcing protocol uses a 50\% duty
cycle (1~ms on, 1~ms off, 5~pulses over 10~ms total), modulated by the
scalar temporal envelope $g(t)$ via the pre-computed Helmholtz projection
(Section~\ref{sec:precomputed_projection}).

Peak displacement reaches $|u_x|^{\max} \approx 3.3$~$\mu$m.
Figure~\ref{fig:tips_displacement_snapshots} shows the displacement
magnitude field at multiple time points, demonstrating shear wave expansion
from the TIPS focal zone.

\begin{figure}[H]
  \centering
  \includegraphics[width=\textwidth]{fig_tips_e_displacement_snapshots.png}
  \caption{Shear wave displacement magnitude snapshots for the TIPS
    configuration.  \textbf{Top row:} $x$--$y$ slices at center elevation.
    \textbf{Bottom row:} $x$--$z$ slices at center lateral position.
    White contours indicate the skull boundary.  The displacement pattern
    reflects the TIPS focal geometry.}
  \label{fig:tips_displacement_snapshots}
\end{figure}

Figure~\ref{fig:tips_displacement_3d} provides a 3D rendering of the
displacement field.

\begin{figure}[H]
  \centering
  \includegraphics[width=0.7\textwidth]{fig_tips_f_displacement_3d.png}
  \caption{3D rendering of the TIPS displacement magnitude.  The
    displacement region is localized in brain tissue behind the skull, with
    the skull rendered as a translucent gray surface.}
  \label{fig:tips_displacement_3d}
\end{figure}

Figure~\ref{fig:tips_displacement_trace} shows the temporal displacement
trace at the peak ARF location, illustrating the pulsed tissue response.

\begin{figure}[H]
  \centering
  \includegraphics[width=0.85\textwidth]{fig_tips_g_displacement_trace.png}
  \caption{Displacement $u_x$ at the peak ARF location for the TIPS
    configuration.  The gray shading indicates the pulsed force envelope
    (50\% duty cycle, 1~ms on / 1~ms off).  Each pulse drives tissue
    displacement followed by viscous relaxation during the off period.}
  \label{fig:tips_displacement_trace}
\end{figure}

\subsection{Comparison with Sparse Phased Array}

The TIPS and sparse array configurations represent two distinct approaches
to transcranial focused ultrasound.  The key differences relevant to shear
wave generation are:

\begin{itemize}[leftmargin=1.5em]
\item \textbf{Aperture geometry.}  The TIPS bowl (92~mm diameter,
  spherical cap, 8~annular elements) produces geometric focusing via the
  curved surface, augmented by electronic delays.  The sparse array
  (1024~point elements on a flat aperture) relies entirely on electronic
  focusing.

\item \textbf{Source pressure.}  The TIPS demonstration uses 400~kPa source
  pressure, compared to 1.5~MPa for the sparse array, reflecting different
  operating regimes (neuromodulation vs.\ clinical ARFI).

\item \textbf{Focal quality.}  The continuous annular aperture of the TIPS
  produces smoother wavefronts through the skull compared to the sparse
  array, which can exhibit grating-lobe-like artifacts from the discrete
  element distribution.

\item \textbf{Skull interaction.}  The TIPS bowl extends further in depth,
  requiring the skull slab to be shifted to avoid geometric intersection.
  This changes the effective skull thickness and angle of incidence relative
  to the sparse array configuration.
\end{itemize}

Despite these differences, the sequential pipeline (acoustic FDTD $\to$
ARF $\to$ shear FDTD) and the shear wave physics remain identical,
demonstrating the generality of the reduced elastodynamic framework for
different transducer configurations.

\subsection{Thermal Neuromodulation Coupling}

The TIPS pipeline is further extended with bioheat-coupled mechanosensitive
neuron simulations, using the \texttt{fullwave.neuromodulation} module to
predict neural firing from the shear wave displacement field.  The heat
source is computed using the same pulse-length and CW envelope normalization
(Eqs.~\ref{eq:isppa_norm}--\ref{eq:cw_correction}) and scaled by the
duty cycle.  A cortical pyramidal neuron model (Pospischil et~al., 2008)
with Piezo1 and TREK-1 mechanosensitive channels predicts spatially
selective firing in brain tissue at the mechanical focus.  Full details of
the neuron model, channel parameterization, and firing results are
presented in the companion neuromodulation report.

\subsection{Acoustic Strain and Strain Gradient}
\label{sec:acoustic_strain}

The preceding sections compute the \emph{elastic} (shear-wave) strain produced
by the slow radiation-force--driven tissue displacement (ms timescale,
$\sim$1--5~m/s shear speed).  A distinct and complementary quantity is the
\emph{acoustic} strain produced directly by the MHz pressure wave itself.  From
linearized acoustics, the volumetric acoustic strain at any point and instant is
related to the local pressure by the equation of state
\begin{equation}
  \varepsilon_{\mathrm{ac}}(\mathbf{x},t)
  = -\frac{p(\mathbf{x},t)}{K(\mathbf{x})},
  \qquad
  K = \rho\,c^2,
  \label{eq:acoustic_strain}
\end{equation}
where $K$ is the bulk modulus.  For soft tissue ($\rho \approx 1000$~kg/m$^3$,
$c \approx 1540$~m/s), $K \approx 2.4$~GPa.  Because $K$ is six orders of
magnitude larger than the shear modulus ($\mu \approx 5$~kPa), the
acoustic strain is extremely small even at therapeutic pressures.  This
relation is instantaneous and exact in the linear regime---no time
averaging or pulse-length normalization is required.

We compute two summary fields from the TIPS pressure output:
\begin{itemize}[leftmargin=1.5em]
  \item \textbf{Peak acoustic strain:} $\varepsilon_{\mathrm{ac}}^{\max}(\mathbf{x}) = \max_t |p(\mathbf{x},t)|/K(\mathbf{x})$,
  \item \textbf{Peak strain gradient magnitude:} $|\nabla\varepsilon_{\mathrm{ac}}|^{\max}(\mathbf{x}) = \max_t |\nabla\varepsilon_{\mathrm{ac}}(\mathbf{x},t)|$,
\end{itemize}
where the spatial gradient $\nabla\varepsilon_{\mathrm{ac}}$ is computed via
central differences on the downsampled grid.  The gradient is taken of the
full strain field $\varepsilon_{\mathrm{ac}} = -p/K$ (rather than just the
pressure gradient) so that spatial variations of the bulk modulus at
tissue interfaces are captured.  To suppress finite-difference artifacts at
the skull--tissue boundary, the skull mask is dilated by one voxel before
zeroing strain and gradient values inside bone.

For the TIPS configuration with 400~kPa source pressure at 1~MHz through the
skull, we obtain:
\begin{itemize}[leftmargin=1.5em]
  \item Peak acoustic strain: $\varepsilon_{\mathrm{ac}}^{\max} \approx 8.8 \times 10^{-4}$
  \item Peak strain gradient: $|\nabla\varepsilon_{\mathrm{ac}}|^{\max} \approx 1.5$~m$^{-1}$
\end{itemize}

\noindent
Figure~\ref{fig:tips_acoustic_strain} shows the peak acoustic strain
distribution.  The strain field mirrors the focal pressure pattern, with the
highest values at the geometric focus and along the beam axis within the skull
aperture.

\begin{figure}[H]
  \centering
  \includegraphics[width=\textwidth]{fig_tips_h_acoustic_strain.png}
  \caption{Peak volumetric acoustic strain $\varepsilon_{\mathrm{ac}}^{\max} = \max_t |p|/K$
    for the TIPS configuration.  \textbf{Left:}~$x$--$y$ slice at center
    elevation.  \textbf{Right:}~$x$--$z$ slice at center lateral position.
    White contours indicate the skull boundary.  The strain pattern
    reflects the acoustic focal geometry.}
  \label{fig:tips_acoustic_strain}
\end{figure}

Figure~\ref{fig:tips_strain_gradient} shows the peak strain gradient
magnitude.  The gradient is largest at the boundaries of the focal zone
and at tissue--bone interfaces where bulk modulus changes abruptly, rather
than at the pressure maximum itself.

\begin{figure}[H]
  \centering
  \includegraphics[width=\textwidth]{fig_tips_i_acoustic_strain_gradient.png}
  \caption{Peak acoustic strain gradient magnitude
    $|\nabla\varepsilon_{\mathrm{ac}}|^{\max}$ for the TIPS configuration.
    \textbf{Left:}~$x$--$y$ slice.  \textbf{Right:}~$x$--$z$ slice.  The
    gradient is concentrated at the edges of the focal zone rather than at
    the pressure peak.}
  \label{fig:tips_strain_gradient}
\end{figure}

\subsection{Acoustic vs.\ Elastic Strain: Comparison and Neuromodulation Mechanisms}
\label{sec:strain_comparison}

Table~\ref{tab:strain_comparison} compares the acoustic (pressure-wave) and
elastic (shear-wave) strain fields computed from the TIPS transcranial
simulation.

\begin{table}[H]
\centering
\caption{Comparison of acoustic and elastic strain magnitudes for the TIPS
  transcranial configuration (400~kPa source, 1~MHz, skull-masked).}
\label{tab:strain_comparison}
\begin{tabular}{lcc}
\hline
Quantity & Acoustic (pressure wave) & Elastic (shear wave) \\
\hline
Peak volumetric strain          & $8.8 \times 10^{-4}$  & $1.8 \times 10^{-4}$ \\
Peak deviatoric strain          & ---                     & $4.4 \times 10^{-4}$ \\
Peak von~Mises strain           & ---                     & $3.6 \times 10^{-4}$ \\
Peak max shear strain           & ---                     & $3.1 \times 10^{-4}$ \\
\hline
\end{tabular}
\end{table}

\noindent
The two strain fields are of the same order of magnitude ($\sim 10^{-4}$),
with the peak acoustic volumetric strain roughly 5$\times$ larger than the
elastic volumetric strain and 2.5$\times$ larger than the elastic von~Mises
strain.  Despite this quantitative similarity, the two fields differ
fundamentally in their spatiotemporal character and the neuromodulation
mechanisms they support.

\subsubsection{Timescale separation}

The \emph{acoustic strain} oscillates at the ultrasound carrier frequency
($f_0 = 1$~MHz, period 1~$\mu$s).  Each acoustic cycle produces a
compressive--rarefactive strain cycle that acts on the membrane for
sub-microsecond durations.  The \emph{elastic strain} evolves on the shear
wave timescale---the shear speed in brain tissue is $c_s \approx 2$~m/s, so
shear wavelengths at the focal zone scale ($\sim$5~mm) correspond to
periods of $\sim$2.5~ms.  The pulsed forcing protocol (1~ms on, 1~ms off)
further shapes the temporal envelope of the elastic strain.  The ratio of
timescales is $\sim$10$^3$, matching the ratio of longitudinal to shear wave
speeds.

\subsubsection{Mechanosensitive ion channels (elastic strain mechanism)}

The shear-wave displacement field acts through the elastic strain tensor to
produce membrane tension (Eq.~\ref{eq:tension}).  Mechanosensitive channels
such as Piezo1 ($T_{1/2} \approx 2.7$~mN/m) and TREK-1 ($T_{1/2} \approx
3.5$~mN/m) gate on the ms timescale, consistent with the shear wave period.
The deviatoric strain ($\sim 4 \times 10^{-4}$) acts through the
area-expansion modulus $K_A$ to produce membrane tension.  For $K_A =
0.25$~N/m, this yields $T \approx K_A \varepsilon_{\mathrm{dev}} \approx
0.11$~mN/m---below the half-activation tensions but in the range where
cumulative or resonant effects may contribute over multiple pulsed
cycles~\cite{kubanek2018neuromodulation}.  This mechanism is intrinsically
\emph{mechanical}: it requires the ARF-driven tissue displacement and is
mediated by the shear modulus of the tissue.

\subsubsection{Flexoelectric coupling (acoustic strain gradient mechanism)}

A distinct hypothesis, developed by Felix et~al.~\cite{felix2022flexo},
proposes that \emph{acoustic strain gradients} directly polarize neural
membranes through the flexoelectric effect.  In this model, the spatial
gradient of the acoustic volumetric strain $\nabla\varepsilon_{\mathrm{ac}}$
produces a polarization density in the lipid bilayer proportional to the
flexoelectric coefficient $\mu_f$:
\begin{equation}
  P = \mu_f \, \nabla\varepsilon_{\mathrm{ac}},
  \label{eq:flexo_polarization}
\end{equation}
where $\mu_f$ has units of C/m and typical values for lipid membranes are
estimated at $\sim$10$^{-18}$--$10^{-17}$~C/m~\cite{felix2022flexo}.  This polarization acts as an effective
transmembrane current that can depolarize or hyperpolarize the membrane,
potentially modulating neural firing without requiring any tissue
displacement or mechanosensitive channel activation.

Several features distinguish this mechanism:
\begin{itemize}[leftmargin=1.5em]
  \item \textbf{Timescale:} The flexoelectric effect operates at the
    acoustic frequency (MHz), orders of magnitude faster than ion channel
    kinetics.  The effective current is the \emph{time-averaged} rectified
    polarization, whose amplitude depends on the strain gradient envelope
    rather than individual acoustic cycles.
  \item \textbf{Spatial pattern:} The strain gradient is largest at the
    \emph{edges} of the focal zone (where the pressure field changes most
    rapidly in space), not at the pressure peak.  This is visible in
    Figure~\ref{fig:tips_strain_gradient} and contrasts with the ARF, which
    peaks at the pressure maximum.
  \item \textbf{Independence from tissue stiffness:} The acoustic strain
    depends on the bulk modulus $K = \rho c^2$ (GPa scale), not the shear
    modulus $\mu$ (kPa scale).  The effect is therefore insensitive to
    variations in tissue viscoelastic properties that strongly modulate the
    shear-wave response.
  \item \textbf{No displacement required:} Because the mechanism operates
    directly on the propagating acoustic wave, it does not require the
    intermediate step of ARF-driven tissue displacement.  Any pressure
    wave---including reflected and scattered components---contributes to the
    strain gradient field.
\end{itemize}

\subsubsection{Comparison figure}

Figure~\ref{fig:tips_mechanism_comparison} provides a side-by-side
comparison of the four spatial fields relevant to neuromodulation:

\begin{figure}[H]
  \centering
  \includegraphics[width=\textwidth]{fig_tips_j_mechanism_comparison.png}
  \caption{Comparison of neuromodulation-relevant fields in the $x$--$y$
    mid-slice.  \textbf{Top left:}~acoustic radiation force (ARF), the
    driving source for shear wave generation.
    \textbf{Top right:}~peak acoustic strain, directly proportional to
    peak pressure.  \textbf{Bottom left:}~peak acoustic strain gradient
    magnitude, relevant for flexoelectric coupling.
    \textbf{Bottom right:}~shear wave displacement magnitude, the end
    result of the ARF$\to$shear wave pipeline.  White contours indicate
    the skull boundary.  Note the qualitatively different spatial patterns:
    ARF and acoustic strain peak at the focal center, while the strain
    gradient is largest at the focal boundary; shear displacement spreads
    over a broader region due to wave propagation.}
  \label{fig:tips_mechanism_comparison}
\end{figure}

\noindent
The four fields have qualitatively different spatial patterns:
\begin{enumerate}[leftmargin=1.5em]
  \item \textbf{ARF and acoustic strain} peak at the geometric focus where
    pressure is highest, producing compact, symmetric distributions.
  \item \textbf{Strain gradient} peaks at the boundary of the focal zone
    where the pressure spatial derivative is largest, forming a ring-like
    or shell-like pattern around the focus.
  \item \textbf{Shear displacement} spreads over a broader region than the
    ARF source due to elastic wave propagation, with the pattern shaped by
    the cumulative history of pulsed forcing and viscous relaxation.
\end{enumerate}

\noindent
These distinct spatial signatures suggest that the two neuromodulation
mechanisms---mechanical (ARF $\to$ shear strain $\to$ membrane tension $\to$
channel gating) and electrical (acoustic strain gradient $\to$ flexoelectric
polarization $\to$ transmembrane current)---would produce different spatial
selectivity profiles.  Experimental measurements of neural activation
patterns during TUS could in principle discriminate between these
mechanisms, since the flexoelectric model predicts peak activation at the
\emph{edges} of the focal zone rather than at the center.

\subsection{Summary}

The TIPS transcranial demo validates the acoustic--elastic--neural pipeline
with a clinically relevant curved-bowl transducer geometry.
The combined pulse-length and CW envelope normalization
(Eqs.~\ref{eq:isppa_norm}--\ref{eq:cw_correction}) provides a
consistent mapping from the transient simulation to CW-equivalent fields for
both radiation force and bioheat.  Peak body force after skull masking is
$\sim$22{,}600~N/m$^3$ and peak shear wave displacement is $\sim$3.3~$\mu$m.

The acoustic strain analysis (Section~\ref{sec:acoustic_strain}) reveals
that the pressure-wave strain ($\varepsilon_{\mathrm{ac}}^{\max} \approx
8.8 \times 10^{-4}$) and the shear-wave elastic strain
($\varepsilon_{\mathrm{dev}}^{\max} \approx 4.4 \times 10^{-4}$) are of the
same order of magnitude but differ by three orders of magnitude in
timescale (MHz vs.\ ms).  The strain gradient field
($|\nabla\varepsilon_{\mathrm{ac}}|^{\max} \approx 1.5$~m$^{-1}$) exhibits
a spatially distinct pattern from the ARF, peaking at the focal boundary
rather than the center.  These results enable quantitative comparison of
the mechanical (ARF--shear--channel gating) and electrical
(flexoelectric--strain gradient) hypotheses for ultrasonic neuromodulation.

The thermal neuromodulation extension couples the displacement field
with bioheat-driven temperature rise and mechanosensitive neuron models,
predicting spatially selective firing volumes localized at the acoustic
focus.

\end{document}
